\section{The actual Type II interface at the ramified coefficient prime}
\label{sec:cubic-269-typeii-interface}

Let \(p=269\), \(K=\mathbb Q(\omega)\),
\(\eta=\eta_K\), \(\chi=\chi_{269}\),
\(j\in\{1,2\}\), and \(\gamma=\psi\kappa_p^e\).
Keep the complete actual cubic ring \(S\) and local prime
\(P_*\) of AC-G65, and every arithmetic, coefficient and
published local hypothesis retained there. In particular
\(e\ge1\), \(p\nmid\operatorname{ord}(\psi)\),
\(\psi_I^2=1\), and \(p-1\nmid2e\).
The original global datum \(D\), when invoked below, is the
literal source \texttt{Dat} for \(K,e,\psi\), not a newly
declared instance for the cubic seed. Source paths have the
pinned root and commit specified in AC-K29-point.

\begin{proposition}[AC-G68-block: middle and determinant at \(269\)]
\label{AC-G68-block}
Let \(f:S\to\overline{\mathbb Q}_p\) be an actual Type II
local point. Write its restriction to \(G_K\) as an endpoint
block \(B_f\) on coordinates \(\{0,2\}\), and a middle
character \(m_f\). Then
\[
                  m_f=1,\qquad \det B_f=\chi^{2j}.
                                                        \tag{RT.1}
\]
Consequently there is an actual continuous integral odd
representation \(V_f:G_{\mathbb Q}\to\operatorname{GL}_2(E)\),
for one finite coefficient extension \(E\), with
\[
 \det V_f=\eta,\qquad
 B_f=\chi^j V_f|_{G_K},\qquad
 V_f(c)=\operatorname{diag}(-1,1).
                                                        \tag{RT.2}
\]
Its restriction to \(G_K\) is irreducible.
\end{proposition}
\begin{proof}
AC-G62-cont, or its constant-dual instance used in AC-G66,
puts all entries in the integers \(O_E\) of one finite
extension, with their prescribed residual coordinates.
In particular the middle character has residual value one.

At the place above \(p\), take the inertia element \(s\)
of (RC.17). It satisfies \(\chi(s)=1\), and
\(a=\gamma(s)\) has \(\overline a^2\ne1\).
The characteristic polynomial supplied by AC-G65-flag is
\((T-a)(T-1)(T-a^{-1})\). Since \(m_f(s)\) reduces to
one, it is exactly one. On the endpoint block the remaining
polynomial is \((T-a)(T-a^{-1})\). Thus the polynomial
projector
\[
 P_1=\frac{(r_f(s)-aI)(r_f(s)-a^{-1}I)}
                 {(1-a)(1-a^{-1})}
             =\operatorname{diag}(0,1,0)
                                                        \tag{RT.3}
\]
is the original middle coordinate projector. All its
denominators are units. In the projector basis of
AC-G65-flag every inertia matrix is upper triangular with
ordered diagonal \(\chi^j\gamma,1,\chi^j\gamma^{-1}\).
Sandwiching it by its middle projector therefore gives
\(P_1r_f(h)P_1=P_1\) for every inertia element \(h\).
In the original block basis the same sandwich is
\(m_f(h)P_1\). Hence \(m_f|_{I_p}=1\).
The determinant polynomial gives
\(\det r_f|_{I_p}=\chi^{2j}|_{I_p}\).
This argument uses the full ramified flag and does not
replace its two endpoint characters by \(\gamma^{\pm1}\).

At any finite place of \(K\) away from \(p\), the actual
residual-kernel killing condition makes inertia factor
through the residual finite image. Its order divides a
prime-to-\(p\) integer, for example
\(3\operatorname{ord}(\psi)\) on \(G_K\).
Both \(m_f\) and \(\det r_f\,\chi^{-2j}\) consequently
have prime-to-\(p\) finite inertia image and residual
value one. Reduction is injective on roots of unity of
such order in \(O_E\), so both characters are trivial
on that inertia as well. They are therefore everywhere
unramified characters of \(G_K\). The ideal class group
of \(K\) is trivial, so class field theory makes both
characters trivial globally. This proves (RT.1).

Use the fixed global complex conjugation \(c\), with
\(r_f(c).\mathrm{mat}=I\), and put
\(C(h)=\chi(h)^{-j}B_f(h)\) for \(h\in G_K\).
The determinant-one calculation of AC-G45 gives
\(C(chc)=D_2C(h)D_2\), where
\(D_2=\operatorname{diag}(-1,1)\). Thus
\(V_f(hc^a)=C(h)D_2^a\) defines the representation
in (RT.2). This uses only unit character values and
integral matrix operations. Type II means that \(B_f\)
is irreducible on \(G_K\); scalar untwisting preserves
this property. All assertions follow.
\end{proof}

\begin{theorem}[AC-G68-point: the old selected point and counting level]
\label{AC-G68-point}
For the point in AC-G68-block, remove \(\chi^j\) from
its endpoint block and keep the middle line \(1\).
Extend this operation to the normalized semidirect
representation as in (GP.5), and call the result
\(\widetilde r_f\). It is an actual point of the
original global source ring: in particular its local
map kills the original selected ideal \(I_t(D)\).
It has the original residual representation and all
seven literal fields of \(D.\texttt{IsPt}\).

Retain additionally the full old
\texttt{AutomorphyUnramifiedCoefficientsAlt}, including
\(e_O=1\), and the original \texttt{CohInp} and
\texttt{GD}. Then \(V_f\) has the original source
\texttt{IsModularDense}, a newform of weight \(2e+1\)
attached after the twist \(\kappa_p^{-e}\), and the
full \texttt{WitAQ}/\texttt{Level} assertions of
\texttt{Adm2A} at level \(N|D_K|\).
No extra factor \(269\) is introduced by removing
the cubic endpoint character.
\end{theorem}
\begin{proof}
\proofstep{The actual old local condition.}
The direct summand \(B_f\) of the representation on
\(G_{K_p}\) is potentially crystalline, with zero
monodromy and finite inertial type
\(\chi^j\tau\oplus\chi^j\tau^{-1}\), where
\(\tau_I=\psi_I\) is the finite character of the
original local datum, literally
\((\texttt{Chr.dirG}(\texttt{comp}\ p\ \psi))\circ D.\iota\)
in \path{UnitaryThreeGroups/RhoKApply.lean},
\texttt{Dat.loc}. After tensoring by \(\chi^{-j}\)
its type is exactly \(\tau\oplus\tau^{-1}\).
This tensor is allowed to be ramified: it removes the
specified finite type, rather than preserving it by an
unramifiedness assertion. Its ordinary flag on
\(G_{K_p}\) has inertia characters
\(\gamma,\gamma^{-1}\), by AC-G65-flag and (RT.1).
Together with the middle line this is the ordered
\(\gamma,1,\gamma^{-1}\) flag of the old local datum.

We verify crystallinity on the exact old subgroup
\(H_L=\ker\eta\cap\ker\tau\).
Let \(L_0\) be its finite fixed field. The untwisted
representation is potentially crystalline on
\(G_{L_0}\) and has trivial inertial type there.
Choose a finite Galois extension \(L_1/L_0\) on which
it is crystalline. Its full crystalline period module
over \(L_{1,0}\otimes E\) has the descent action of
\(\operatorname{Gal}(L_1/L_0)\). The inertia subgroup
acts trivially, because its action on this module is
precisely the restricted potentially crystalline
inertial type. The remaining unramified quotient acts
semilinearly. Ordinary finite Galois descent gives a
full-rank module over \(L_{0,0}\otimes E\).
Its invariants are exactly
\(D_{\mathrm{cris},L_0}\), since they are the period
vectors invariant under \(G_{L_0}\). Thus the
representation is crystalline already on \(H_L\).
This is the period descent also used in
AC-G57-TypeII; no assumption that \(\chi\) itself
is crystalline has been made.

The untwisted local semidirect endpoint determinant is
one on every element, and the middle line is one.
Its residue is the old \(r_0(\eta,\gamma)\).
The source \texttt{GD.ine} and \texttt{GD.gen} imply
\[
           \overline\gamma^{\,2}|_{I_p}
                =\overline\kappa_p^{\,2e}|_{I_p}\ne1.
                                                        \tag{RT.4}
\]
All hypotheses of AC-G48-TypeII now hold. Apply that
proposition directly, rather than its separate
unramified-cubic corollary: the original local point
kills \(I_t(D)\). More explicitly, its saturated
ordinary endpoint line gives the strict basis in
(G48.7); the diagonal unramified retagging leaves
\(I_Y\) unchanged, and the crystalline extension
contraction kills the selected prime through the
source separator \(yI_t\subset I_Y\).
Finite coefficient descent there proves membership
in the original ideal over \(O\).
We do not identify \(I_t\) with all of \(I_Y\),
or assert that an arbitrary point on \(I_t\) kills
the unsaturated ordinary equation ideal by definition.

\proofstep{All global point fields and the actual ring map.}
The complete proof of AC-K29-kill applies at \(269\):
it requires only \(p\ne3\), and its cubing argument
concerns inertia away from the coefficient prime.
It proves the old residual-kernel killing condition,
strict old residual congruence, and normalization at
every old marked element. Continuity, integrality and
the multiplier were preserved by the semidirect twist.
The just proved local membership supplies exactly
(GP.11), the seventh field of \texttt{IsPt} in
\path{UnitaryThreeGroups/RhoKApply.lean}.

For precision the coefficient comparison is made in
\(A=O+\mathfrak m_E\subset O_E\). The full argument
of AC-K29-point makes \(A\) a finite complete local
\(O\)-algebra with the original residue field, puts
the old local map and every untwisted matrix in \(A\),
and applies the literal \texttt{GLift}/\texttt{BoxCond}
universal property. It gives an actual local map to
the old global ring, and its entry-algebra restriction
is the \texttt{IsPush} point required by
\texttt{LemWAlt}. Enlarging the value field of a point
has not changed the original coefficient datum or
removed its \(e_O=1\) hypothesis.

\proofstep{The unchanged old modularity and level theorem.}
Under the additional source input in the statement,
\texttt{AutomorphyUnramifiedCoefficientsAlt} supplies
the old marked element and \texttt{LemWAlt}.
The new actual point satisfies its hypotheses by the
preceding two steps. Thus it gives
\texttt{IsModularDense} for every old \texttt{IsExt},
including the explicit \(V_f\) in (RT.2).
Apply the same complete
\texttt{adm2B\_rhoK\_of\_typeII\_of\_modular\_extensions}
proof in \path{ModularityArgument/ThresholdArguments/DenseThresholdCases/AdmissibleOfTypeII.lean}
used in AC-K29-typeII. It checks the actual ordinary
flag, away inertia and determinant, and yields the
stated newform, \texttt{WitAQ} and original level
\(N|D_K|\). Its application concerns the untwisted
old point. The extra cubic conductor is therefore
absent; any factor \(269\) already in the original
\(N\) for \(\psi\) remains exactly as specified.
\end{proof}

\begin{theorem}[AC-G68: the exact remaining Type II rigidity coefficient]
\label{AC-G68}
Retain all assumptions of AC-G68-point, including
the old automorphy input, \(e_O=1\), and additionally
the literal old \texttt{RigiditySecondCase}.
For every actual Type II point \(f\), there is a
natural isomorphism, for finite coefficient fields
containing the point and character values,
\[
 H_f^1(\mathbb Q,\operatorname{ad}r_f)
       \simeq H_f^1(K,V_f|_{G_K}\otimes\chi^j).
                                                        \tag{RT.5}
\]
Here the local conditions are the actual source
\texttt{Af.IsBK} conditions, including their
\(B_{\mathrm{cris}}\) kernel at \(269\).

Allow finite coefficient extensions when taking actual
lifts, and take the direct limit of the right side of
(RT.5) over those extensions. Let \(\mathcal T_f\)
be the image in that direct limit of all actual algebraic dual lifts
\(F:S\to\overline{\mathbb Q}_p[\epsilon]/\epsilon^2\)
with first component \(f\). Then
\[
 \texttt{Rigid}(r^{\rm univ},f)
                  \quad\Longleftrightarrow\quad
                         \mathcal T_f=0.              \tag{RT.6}
\]
In particular vanishing of the entire right side of
(RT.5) suffices, but no surjectivity of the actual
tangent image onto it is asserted. These statements
are pointwise for every coefficient extension,
marked element and ring/model choice admitted by
AC-G65. They assert no mixed vanishing at \(269\).
\end{theorem}
\begin{proof}
The three coordinate projections of (GT.3)--(GT.4)
give an actual equivariant direct sum
\[
 \operatorname{ad}r_f
       =\mathcal A_f\oplus\mathcal C_f\oplus\mathcal R_f,
 \quad \mathcal C_f=E(\eta),\quad
 \mathcal R_f\simeq
       \operatorname{Ind}_{G_K}^{G_{\mathbb Q}}
                  (V_f|_{G_K}\otimes\chi^j).
                                                        \tag{RT.7}
\]
The last identity is the full coset calculation of
AC-G45-cohomology. It does not require \(\chi\) to
be unramified at the coefficient prime. Its finite
coset cochain map commutes with tensoring by
\(B_{\mathrm{cris}}\), so it preserves the required
local kernels also at this ramified prime.

The endpoint scalar \(\chi^j\) cancels in the
corner conjugation, and the inverse-transpose
action on the other coset is unchanged. Hence
\(\mathcal A_f\) is exactly the corner adjoint
of the old point \(\widetilde r_f\).
Every premise of the old
\path{ModularityArgument/ThresholdArguments/DenseThresholdCases/RigEnd3.lean},
\texttt{RigiditySecondCase}, has now been supplied:
the literal old datum and ring, actual local point,
Type II, the continuous \texttt{IsExt} representation
\(V_f\), its determinant \(\eta\), and its actual
\texttt{IsModularDense}. Its \texttt{AdZero2}
conclusion applies through exactly the corner
maps \texttt{pA}, \texttt{sA} and the symmetric
pairing \(J_2\) in the full proof of
\path{ModularityArgument/ThresholdArguments/DenseThresholdCases/TypeIIOfBlochKatoLocalPoint.lean}.
Consequently \(H_f^1(\mathbb Q,\mathcal A_f)=0\).
This uses its corner step only; the old
\texttt{kUse'} assertion for the two rational
character twists is not applied to \(\mathcal R_f\).
Also AC-G53-Artin with weight zero gives
\(H_f^1(\mathbb Q,E(\eta))=0\).
If those source conclusions are initially expressed
over \(\overline{\mathbb Q}_p\), a cobounding
matrix has entries in a finite extension; applying
the normalized coefficient-field trace descends
the coboundary to the original finite field.
Thus the stated finite-coefficient assertions hold.

All three projections and inclusions preserve
cocycles and local coboundaries, including the
period-ring cobounding tensors. The Selmer group
therefore respects (RT.7), and its first two
summands vanish. Shapiro with the specified local
kernels proves (RT.5).

For every actual dual lift, AC-G66-BK supplies
continuity and precisely those local conditions.
Its corner and middle cocycles are coboundaries
by the preceding argument. If its mixed class
is zero, sum the three cobounding matrices to
obtain the full coboundary, with the source sign
\(z(g)=X-\operatorname{ad}(r_f)(g)X\).
Conversely the mixed projection of any full
coboundary is a coboundary. Quantifying over
every actual lift proves (RT.6), exactly as in
AC-K29-tangent, now with its local-condition,
corner and middle premises established for
the actual \(269\)-adic ring.
\end{proof}

\paragraph{Arithmetic and geometric scope.}
The conditional phrase concerning the actual Type II
point in AC-G65 has been removed: its middle and
determinant identities, old selected component,
old global point and original modular counting level
are proved under the retained old inputs.
The remaining \(269\)-adic rigidity assertion is
the exact mixed image (RT.6), or the sufficient
rank-two-over-\(K\) vanishing in (RT.5).
The ordinary CM Euler constructions at the other
coefficient primes do not prove that assertion:
the ramified theta factor at \(269\) remains outside
their verified coefficient specialization.
Neither Type I Hecke containment, \texttt{M2}, nor
a uniform integral index estimate follows here.
Geometrically, endpoint untwisting tensors a
decomposed point's rank-two lisse coefficient by
the inverse cubic rank-one system, and the selected
component comparison is a statement about its
actual formal deformation point. The previously
audited SGA~1 finite-cover dictionary supplies no
additional Selmer vanishing or automorphy theorem.
